\documentclass[11pt,reqno]{amsart}

\usepackage[a4paper,margin=28mm]{geometry}
\usepackage{amsmath,amssymb,mathtools,mathrsfs}
\usepackage{enumitem}
\usepackage{microtype}
\usepackage{xurl}
\usepackage[colorlinks=true,linkcolor=blue,citecolor=blue,urlcolor=blue]{hyperref}

\numberwithin{equation}{section}
\setlist[enumerate]{label=\textup{(\roman*)},leftmargin=2.2em}

\newtheorem{theorem}{Theorem}[section]
\newtheorem{proposition}[theorem]{Proposition}
\newtheorem{lemma}[theorem]{Lemma}
\newtheorem{corollary}[theorem]{Corollary}
\newtheorem{example}[theorem]{Example}
\newtheorem{remark}[theorem]{Remark}
\theoremstyle{definition}
\newtheorem{definition}[theorem]{Definition}

\newcommand{\N}{\mathbb N}
\newcommand{\Nzero}{\mathbb N_0}
\newcommand{\C}{\mathcal C}
\newcommand{\D}{\mathcal D}
\newcommand{\Acal}{\mathcal A}
\newcommand{\Lcal}{\mathcal L}
\newcommand{\Patt}{\mathsf P}
\newcommand{\hseq}{h_{\mathrm{seq}}}
\newcommand{\hpat}{h_{\mathrm{pat}}^{*}}
\newcommand{\hinv}{h_{\mathrm{inv}}}
\newcommand{\htop}{h_{\mathrm{top}}}
\newcommand{\Bor}{\mathrm{Bor}}
\newcommand{\clop}{\mathrm{clop}}
\newcommand{\lang}{\mathrm{lang}}
\newcommand{\supp}{\operatorname{supp}}
\newcommand{\doi}[1]{\href{https://doi.org/#1}{doi:\nolinkurl{#1}}}

\title{Sparse-Time Complexity of Invariant Control Languages}
\author{Ziqing Ding}
\address{School of Mathematical Sciences, Nanjing Normal University,
Nanjing 210023, China}
\date{}
\subjclass[2020]{93C55, 37B10, 37B40, 94A17}
\keywords{invariance entropy, invariant partition, sequence entropy,
maximal pattern complexity, symbolic control, minimax equality}

\begin{document}

\begin{abstract}
We study sparse observations of admissible name languages associated with
invariant partitions of deterministic control systems. For every countable
family of nonempty finite-alphabet pattern systems, one increasing coordinate
sequence simultaneously attains their maximal-pattern entropy rates. A
coordinate-deletion argument removes the need for shift invariance. This
yields a minimax identity at a fixed sampling period for any nonempty class
of invariant partitions with countably many languages up to relabeling.
We also realize every one-sided subshift as a compact matching-control
system. Kawan's exact-containment invariance entropy of this system equals
the topological entropy of the subshift, and its period-one sequence and maximal-pattern infima over all
Borel invariant partitions equal the corresponding coordinate rates.
The Thue--Morse subshift consequently gives zero invariance entropy and
period-one pattern infimum $\log 2$. We give a direct binary construction of
an attaining sparse schedule. Finally, allowing all sampling periods makes
the unrestricted Borel pattern infimum collapse to Kawan's invariance
entropy, which explains the role of the fixed-period formulation.
\end{abstract}

\maketitle

\section{Introduction}\label{sec:intro}

Ordinary entropy measures exponential complexity along consecutive times.
A zero value does not imply that observations made far apart are simple:
a system may have subexponentially many long orbit names and still realize
exponentially many patterns on well-chosen sets of observation times.
Sequence entropy fixes one increasing sequence of times, whereas maximal
pattern entropy optimizes over finite sets of a prescribed cardinality.
The normalization is by the number of observations. Thus these notions
examine information that an elapsed-time entropy rate can miss.

The topological sequence entropy of Goodman \cite{Goodman1974} and the
maximal pattern complexity of Kamae and Zamboni
\cite{KamaeZamboniWords2002,KamaeZamboniSystems2002} provide the dynamical
background. Huang and Ye \cite[Theorem~2.1]{HuangYe2009} proved that, for
each fixed finite open cover of a compact invertible system, maximal
pattern entropy is an attained supremum of sequence entropies. They also
proved its measurable-partition counterpart. Their paper develops much
more than this attainment statement: its combinatorial lemmas relate
pattern growth to independence and to intrinsic sequence-entropy tuples.
The input relevant here is precisely the attained-supremum theorem and
its long-block construction.

The measure-theoretic study of sequence entropy and subsequence generators
includes Hulse's earlier work \cite{Hulse1982}. A measurable attainment
theorem is given by Huang, Maass and Ye
\cite[Theorem~2.3]{HuangMaassYe2004}. The existence of a common
sequence, by itself, is not a novelty claim of this paper. Nor are arbitrary
sets of finite-alphabet sequences new objects: such pattern sets occur,
for example, in Kamae's study of uniform complexity \cite{KamaeUniform2009}.
Our combinatorial statement concerns the individual maximal counting rates
of an arbitrary countable family of such sets, without a common ambient
transformation or mixing assumptions.

The control-theoretic question is different from fixing an observation
partition of one dynamical system. For a controlled-invariant compact set
$Q$, a finite invariant partition assigns a safe control block to each
atom. Iterating this rule generates admissible atom names, while retaining
safety at all intervening times. Kawan's invariant-cover framework recovers
his exact-containment invariance entropy by minimizing consecutive-name entropy over
Borel invariant partitions and sampling periods
\cite[Definitions~2.8--2.9 and Theorem~2.3]{Kawan2013}.
In the terminology of Yao \cite[Section~2]{Yao2026}, this exact-containment
quantity is called \emph{strict invariance entropy}; the neighborhood-based
\emph{ordinary invariance entropy} there is Kawan's \emph{outer invariance
entropy} \cite[Definition~2.3]{Kawan2013}. We retain Kawan's notation
$h_{\mathrm{inv}}$ throughout.
We ask what happens when the same admissible languages are observed only
on selected block coordinates.

Passing from a fixed strategy to an infimum over strategies creates a
quantifier problem. A maximizing observation sequence for one invariant
partition may be unsuitable for another. An identity of the form
$\sup_A\inf_{\mathcal C}=\inf_{\mathcal C}\sup_A$ therefore needs a common
attainment argument. There is a second issue: the natural name systems
are one-sided and their shift maps need not be onto. Translating a
maximizing coordinate set to later times can then lose information.
The coordinate-deletion estimate below handles this issue without a
surjectivity assumption.

The results form the following chain.
\begin{enumerate}[label=\textup{(\Alph*)}]
\item \emph{Simultaneous entropy-rate attainment.}
For nonempty sets $\Omega_j\subset E_j^{\mathbb N_0}$ over finite
alphabets and a fixed scale $\tau>0$, there exists one increasing sequence
$A$ such that
\begin{align*}
 h_{\mathrm{seq}}^A(\Omega_j;\tau)
 =h_{\mathrm{pat}}^*(\Omega_j;\tau)\qquad(j\ge1).
\end{align*}
The proof revisits each member infinitely often. At each stage it deletes
the already passed coordinates of a maximizing finite set and chooses
its size large enough to dilute the previous observations.

\item \emph{A fixed-period minimax identity.}
For a nonempty class $\mathfrak C$ of period-$\tau$ Borel invariant
partitions having at most countably many admissible languages up to
alphabet relabeling,
\begin{align*}
 \max_A\inf_{\mathcal C\in\mathfrak C}
 h_{\mathrm{seq}}^A(\mathcal C;Q)
 =\inf_{\mathcal C\in\mathfrak C}h_{\mathrm{pat}}^*(\mathcal C;Q).
\end{align*}
The maximum is attained by a sequence which attains every individual
partition rate. A compact metric state constraint and a countable control
alphabet give a useful sufficient condition for the clopen class.

\item \emph{Exact control realization.}
Every nonempty one-sided subshift $Y$ over a finite alphabet admits a
compact matching-control realization with
$r_{\mathrm{inv}}(N,Y)=\#\mathcal L_N(Y)$ and
$h_{\mathrm{inv}}(Y)=h_{\mathrm{top}}(Y)$.
For every period-one Borel invariant partition, its finite-prefix name
system factors onto $Y$ through the assigned control labels. The
first-coordinate cylinder partition attains the resulting lower bounds
for all finite pattern counts. This computes the Borel and clopen
sequence and pattern infima, without a countability hypothesis.

\item \emph{Strict separation at period one.}
For the Thue--Morse realization,
\begin{align*}
 0=h_{\mathrm{inv}}(Y)
 <\log2=\inf_{\mathcal C\in\mathfrak C_{\mathrm{Bor}}(Y,1)}
 h_{\mathrm{pat}}^*(\mathcal C;Y).
\end{align*}
The explicit schedule $a_i=4^{i-1}-1$ realizes all binary patterns on
every initial segment. Both the zero consecutive-time rate and this
sparse-time calculation are proved directly.
\end{enumerate}

The elementary symbolic matching mechanism should not itself be regarded
as new. For comparison, Yao \cite[Theorem~3.1 and Proposition~4.5]{Yao2026}
uses exact symbolic agreement in a continuous-time construction separating
strict safety from neighborhood-based invariance. Our result concerns a
different optimization: the period-one name languages of \emph{every}
Borel invariant partition and their exact projection onto an arbitrary
one-sided subshift. The Thue--Morse symbolic input is classical; its
role here is to give a fully computed example of the control-language
infimum.

The contribution is thus a simultaneous counting argument coupled with an
exact realization of sparse invariant-control languages. The proofs are
elementary extensions of established entropy methods. We do not claim an
unrestricted Borel minimax theorem for general control systems. We also
do not identify a fixed-period sparse rate with a minimum communication
rate per elapsed time. In fact, Proposition~\ref{prop:all-period-collapse} shows that
minimizing the pattern rate over all periods recovers Kawan's invariance
entropy. Fixing the sampling period is essential for the separation above.

Section~\ref{sec:background} recalls the entropy background.
Section~\ref{sec:control} constructs the admissible name systems.
Sections~\ref{sec:comb} and \ref{sec:minimax} prove (A) and (B).
Sections~\ref{sec:realization} and \ref{sec:thue-morse} prove (C) and (D).
Section~\ref{sec:comparison} explains the all-period collapse and the
remaining limitations.

\section{Entropy background and the control interface}\label{sec:background}

Write $\mathbb N=\{1,2,\ldots\}$ and $\mathbb N_0=\{0,1,2,\ldots\}$.
All logarithms in this paper are natural. For a finite open cover
$\mathcal V$ of a compact space, let $N(\mathcal V)$ be the minimum
cardinality of a subcover. The join of finitely many covers consists of
all intersections of one member from each cover, with empty intersections
removed. For a continuous map $T:Z\to Z$ and a finite open cover
$\mathcal V$, ordinary cover entropy is
\begin{align*}
 h_{\mathrm{top}}(T,\mathcal V)
 =\lim_{n\to\infty}\frac1n\log N\left(
 \bigvee_{i=0}^{n-1}T^{-i}\mathcal V\right).
\end{align*}
Submultiplicativity of cover numbers gives this limit. Taking the
supremum over finite open covers defines $h_{\mathrm{top}}(T)$.
For $A=(a_i)_{i\ge1}$ with $0\le a_1<a_2<\cdots$, set
\begin{align*}
 h_{\mathrm{top}}^A(T,\mathcal V)
 &=\limsup_{n\to\infty}\frac1n\log N\left(
     \bigvee_{i=1}^nT^{-a_i}\mathcal V\right),\\
 p_{Z,\mathcal V}^*(n)
 &=\max_{\substack{F\subset\mathbb N_0\\\#F=n}}
    N\left(\bigvee_{i\in F}T^{-i}\mathcal V\right),\\
 h_{\mathrm{top}}^*(T,\mathcal V)
 &=\lim_{n\to\infty}\frac1n\log p_{Z,\mathcal V}^*(n).
\end{align*}
The sequence in the last line is subadditive before division by $n$;
the arbitrary sequence in the first line need not have that property.
This is why a limit and a limsup appear in different definitions.

Huang--Ye's Theorem~2.1(1) states, under their convention that $T$ is a
homeomorphism of a compact metric space, that
\begin{equation}\label{eq:HY-background}
 h_{\mathrm{top}}^*(T,\mathcal V)
 =\max_A h_{\mathrm{top}}^A(T,\mathcal V).
\end{equation}
Its proof selects long nearly maximizing time sets and separates them
by translations. Surjectivity preserves the cover number under a common
translation. The previous sets contribute a vanishing fraction of the
selected observations. For a non-surjective map the translation step
need not preserve cover numbers; our proof uses deletion instead.

If $\mathcal V$ is a clopen partition, the join count is exactly the
number of names seen on the selected coordinates. An arbitrary open
cover may have overlaps, so its minimum subcover count must not be
replaced by the number of nonempty intersections. Similarly, the
measurable version of \eqref{eq:HY-background} uses Shannon entropy
$H_\mu(\alpha)=-\sum_{D\in\alpha}\mu(D)\log\mu(D)$, not the logarithm
of a support cardinality. Our language counts use the latter quantity.
Theorem~\ref{thm:A} is a counting statement, not a measurable
common-sequence theorem.

For a finite-alphabet subshift $Y$, the first-coordinate cylinders give
one particular clopen partition. Its coordinate sequence rate is a
cover-level quantity. Unlike ordinary topological entropy, sequence
entropy need not be computed by a generating partition: replacing
one-coordinate observations by fixed-length windows can change the
normalized sparse rate. We therefore keep the subscript
``$\mathrm{coord}$'' in the realization theorem and do not silently
identify this rate with the supremum over all open covers.

Kawan's Chapter~2 supplies the other interface. Definition~2.2 defines
invariance entropy from finite-horizon spanning controls.
Definitions~2.8 and 2.9 introduce invariant covers and their admissible
word entropies. Proposition~2.21 proves word-count submultiplicativity;
Proposition~2.22 bounds invariance entropy by invariant-cover entropy.
Lemma~2.3 disjointizes a cover without increasing its name entropy.
It does \emph{not} turn arbitrary nonmeasurable sets into Borel sets.
Borel measurability in Theorem~2.3 comes from the safe sets associated
with finitely many spanning controls and their successive differences.
The theorem takes an infimum over sampling periods as well as
partitions. We provide the discrete-time safe-set argument in
Section~\ref{sec:comparison}, where its application to pattern rates is
also made explicit.

\section{Control systems and invariant name languages}\label{sec:control}

Kawan's invariance entropy and consecutive-name partition entropy are measured in nats
per physical time unit; the sparse-time normalizations are specified
separately below.  Let $X$ be a
metrizable space, let $U$ be a nonempty control set, and suppose that
$f_u:=f(\,\cdot\,,u):X\to X$ is continuous for every $u\in U$.  For a finite
control word $\omega=(\omega_0,\ldots,\omega_{N-1})\in U^N$, set
\begin{align*}
 \varphi(0,x,\omega)=x,
 \qquad
 \varphi(n+1,x,\omega)
 =f\bigl(\varphi(n,x,\omega),\omega_n\bigr)
 \quad(0\le n<N).
\end{align*}
A nonempty compact set $Q\subset X$ is \emph{controlled invariant} if every
$x\in Q$ admits an infinite control whose complete forward trajectory stays in
$Q$.

For $N\in\N$, let $r_{\rm inv}(N,Q)$ be the least cardinality of a set
$S\subset U^N$ such that for every $x\in Q$ some $\omega\in S$ satisfies
\begin{align*}
 \varphi(n,x,\omega)\in Q\qquad(0\le n\le N).
\end{align*}
If no finite such set exists, put $r_{\rm inv}(N,Q)=+\infty$.  The
discrete-time invariance entropy is
\begin{equation}\label{eq:hinv-definition}
 \hinv(Q)=\limsup_{N\to\infty}\frac1N\log r_{\rm inv}(N,Q).
\end{equation}
Throughout, the admissible control space is the full sequence space
$\mathcal U=U^{\Nzero}$; equivalently one may use Kawan's two-sided full
control space and retain its nonnegative coordinates.  Every finite word
therefore extends to an admissible infinite control, and finite-horizon
dependence identifies the preceding spanning number with Kawan's.  Thus
\eqref{eq:hinv-definition} is Kawan's definition with the logarithm converted
from base two to base $e$; see \cite[Definition~2.2, p.~44]{Kawan2013}.

Fix a sampling period $\tau\in\N$.  A \emph{finite invariant cover} of $Q$
with period $\tau$ is a triple $\C=(\Acal,\tau,\nu)$, where $\Acal$ is a finite
cover of $Q$ and $\nu(D)\in U^\tau$ is assigned to each $D\in\Acal$, such that
\begin{equation}\label{eq:invariant-cover}
 \varphi(s,x,\nu(D))\in Q
 \qquad(x\in D,\ 0\le s\le\tau).
\end{equation}
This is the discrete-time specialization of Kawan's Definition~2.8
\cite[p.~79]{Kawan2013}.  An \emph{invariant partition} is an invariant cover
whose underlying family is a finite Borel partition of $Q$ into nonempty
atoms.  It is \emph{clopen} if its atoms are clopen in $Q$.  We write
\begin{align*}
 \mathfrak C_{\Bor}(Q,\tau)&=\bigl\{\text{finite Borel invariant partitions of
 }Q\text{ of period }\tau\bigr\},\\
 \mathfrak C_{\clop}(Q,\tau)&=\bigl\{\C\in\mathfrak C_{\Bor}(Q,\tau):
 \C\text{ is clopen}\bigr\},\\
 \mathfrak C_{\Bor}(Q)&=\bigcup_{\tau\in\N}\mathfrak C_{\Bor}(Q,\tau).
\end{align*}
We use the extended-value convention $\inf\varnothing=+\infty$.

Let $\C=(\{D_i:i\in I_\C\},\tau,\nu)$ be an invariant partition, with
$I_\C$ a finite nonempty alphabet.  For
$a=(a_0,\ldots,a_{N-1})\in I_\C^N$, concatenate the assigned blocks in that
order.  The word $a$ is \emph{admissible} if some $x\in D_{a_0}$ follows this
concatenated control and satisfies
\begin{equation}\label{eq:admissible-word}
 \varphi(j\tau,x,\nu(D_{a_0})\cdots\nu(D_{a_{N-1}}))
 \in D_{a_j}\qquad(0\le j<N).
\end{equation}
The set of such words is $W_N(\C;Q)$.  These are the partition words of
Kawan's Definition~2.9 \cite[pp.~79--80]{Kawan2013}.

\begin{lemma}[Safety, suffixes, and right extension]\label{lem:safety}
Let $\C$ be an invariant partition.
\begin{enumerate}
\item A witness for $a\in W_N(\C;Q)$ remains in $Q$ at every physical time
from $0$ through $N\tau$ under the concatenated blocks.
\item Every prefix and every suffix of an admissible word is admissible.
\item Every $a\in W_N(\C;Q)$ has a right extension in
$W_{N+1}(\C;Q)$.  In particular, every finite admissible word has an
infinite right extension.
\item For $M,N\in\N$,
\begin{equation}\label{eq:word-submult}
 \#W_{M+N}(\C;Q)
 \le \#W_M(\C;Q)\,\#W_N(\C;Q).
\end{equation}
\end{enumerate}
\end{lemma}

\begin{proof}
On block $j$, condition \eqref{eq:invariant-cover}, applied to the state in
$D_{a_j}$, keeps the trajectory in $Q$ throughout the next $\tau$ steps.
Induction proves (i).  Prefixes are immediate.  A suffix beginning at block
$j$ is witnessed by the state at time $j\tau$, which proves (ii).

At time $N\tau$ the terminal state belongs to a unique atom $D_b$.
Appending $b$ and its assigned block gives an admissible word of length
$N+1$, proving (iii) without compactness.  Finally, an admissible word of
length $M+N$ is determined injectively by its length-$M$ prefix and
length-$N$ suffix.  Part (ii) places these in $W_M$ and $W_N$, respectively,
and gives \eqref{eq:word-submult}.
\end{proof}

Define the finite-prefix name system
\begin{equation}\label{eq:name-system}
 \Sigma_\C=\bigl\{a\in I_\C^{\Nzero}:a|_{[0,N-1]}\in W_N(\C;Q)
 \text{ for every }N\in\N\bigr\}.
\end{equation}
We give $I_\C$ the discrete topology and $\Sigma_\C$ the subspace topology
from $I_\C^{\Nzero}$.  Its complement is a union of cylinders determined by
forbidden finite prefixes, so $\Sigma_\C$ is closed.
Lemma~\ref{lem:safety}(iii) shows that every word in $W_N$ is the initial
word of some element of $\Sigma_\C$.  Lemma~\ref{lem:safety}(ii) shows
$\sigma(\Sigma_\C)\subset\Sigma_\C$; surjectivity need not hold.  For Borel
atoms, an element of $\Sigma_\C$ need not have one common state witness for
all its prefixes.  Our definitions use only finite projections and do not
assume such a witness.

There is a useful concrete description of $\Sigma_\C$. The feedback
block map and its consecutive atom names also appear in the numerical
framework of \cite[Section~2]{TomarKawanZamani2022}. Define
\begin{align*}
 T_\C(x)=\varphi(\tau,x,\nu(D_i))\quad\text{if }x\in D_i,
\end{align*}
and its atom itinerary $\iota_\C(x)$ by
$(\iota_\C(x))_j=i$ if $T_\C^j(x)\in D_i$.
Each branch of $T_\C$ is continuous, and $T_\C$ is Borel measurable;
it need not be continuous when the atoms are merely Borel. Every finite
admissible word is the prefix of the actual itinerary of its witness.
Conversely every actual itinerary has admissible prefixes. It follows that
\begin{equation}\label{eq:itinerary-closure}
 \Sigma_\C=\overline{\iota_\C(Q)}.
\end{equation}
Indeed, each prefix cylinder meeting $\Sigma_\C$ meets
$\iota_\C(Q)$, which proves density, and all actual itineraries belong
to the closed set $\Sigma_\C$.

\begin{remark}[Infinite witnesses]\label{rem:witness}
If the atoms are clopen, $T_\C$ and $\iota_\C$ are continuous.
Compactness of $Q$ makes $\iota_\C(Q)$ closed, so every element of
$\Sigma_\C$ has a common state witness. Equivalently, the sets of witnesses
for its finite prefixes are nested nonempty compact sets.

This conclusion can fail for a Borel partition even in a compact
continuous system with one control. Take $Q=X=[0,1]$,
$f(x)=\min\{2x,1\}$, $D_0=(0,1/2)$ and $D_1=Q\setminus D_0$.
Both atoms receive the unique control, so the partition is invariant.
For every $N$, the point $x_N=2^{-N-1}$ witnesses the word $0^N$:
for $0\le j<N$, $f^j(x_N)=2^{j-N-1}\in(0,1/2)$.
Thus $0^\infty\in\Sigma_\C$. No point witnesses this infinite name:
zero is outside $D_0$, and every positive point eventually leaves $D_0$.
\end{remark}

For a nonempty finite $F\subset\Nzero$, put
\begin{equation}\label{eq:control-pattern-count}
 \Patt_\C(F)=\#\{a|_F:a\in\Sigma_\C\}.
\end{equation}
Equivalently, if $m=\max F$, this is the number of projections of
$W_{m+1}(\C;Q)$ onto $F$.  Thus a counted pattern always has a finite witness
that remains in $Q$ also between the selected block times.

Let
\begin{align*}
 \mathscr S=\{A=(a_n)_{n\ge1}:0\le a_1<a_2<\cdots, a_n\in\Nzero\}.
\end{align*}
Define
\begin{align}
 p_\C^*(n)&=\max_{\substack{F\subset\Nzero\\\#F=n}}\Patt_\C(F),
 \label{eq:pstar-control}\\
 \hpat(\C;Q)&=\lim_{n\to\infty}\frac1{n\tau}\log p_\C^*(n),
 \label{eq:hpat-control}\\
 \hseq^A(\C;Q)&=\limsup_{n\to\infty}\frac1{n\tau}
 \log\Patt_\C(\{a_1,\ldots,a_n\}).
 \label{eq:hseq-control}
\end{align}
The existence of the limit in \eqref{eq:hpat-control} follows from
Lemma~\ref{lem:pattern-submult} below.  Although there are infinitely many
$n$-element subsets of $\Nzero$, the maximum in \eqref{eq:pstar-control}
exists: its possible values form a nonempty subset of
$\{1,\ldots,(\#I_\C)^n\}$.

The entries of $A$ are \emph{block indices}; the corresponding physical
observation times are $\tau A=(\tau a_n)_{n\ge1}$.  The normalization in
\eqref{eq:hseq-control} is $n\tau$, not the elapsed horizon $a_n\tau$.
For fixed $\tau$, $A\mapsto\tau A$ is a bijection between block-index
schedules and physical schedules contained in $\tau\Nzero$.  If sampling
periods vary, the same block-index sequence represents different physical
schedules.

The ordinary consecutive-name entropy is
\begin{equation}\label{eq:ordinary-partition-entropy}
 h(\C;Q)=\lim_{N\to\infty}\frac1{N\tau}\log\#W_N(\C;Q),
\end{equation}
whose limit exists by \eqref{eq:word-submult} and Fekete's lemma.

\section{Simultaneous attainment for pattern systems}\label{sec:comb}

A \emph{finite-alphabet pattern system} is a nonempty set
$\Omega\subset E^{\Nzero}$ with $E$ finite and nonempty.  No topology or
shift invariance is imposed.  For finite $F\subset\Nzero$, write
\begin{align*}
 \Patt_\Omega(F)=\#\{x|_F:x\in\Omega\},
 \qquad
 p_\Omega^*(n)=
 \max_{\substack{F\subset\Nzero\\\#F=n}}\Patt_\Omega(F).
\end{align*}
With a fixed scale $\tau>0$, define
\begin{align*}
 \hpat(\Omega;\tau)&=\lim_{n\to\infty}
       \frac1{n\tau}\log p_\Omega^*(n),\\
 \hseq^A(\Omega;\tau)&=\limsup_{n\to\infty}
       \frac1{n\tau}\log
       \Patt_\Omega(\{a_1,\ldots,a_n\}).
\end{align*}
The maximum exists: its possible values form a nonempty subset of the finite
integer set $\{1,\ldots,(\#E)^n\}$.

We use $\Patt_\Omega(\varnothing)=1$ when needed.
Taking the product-topology closure of $\Omega$ changes none of its
finite projections: a cylinder meeting $\overline\Omega$ already meets
$\Omega$. Thus dispensing with closedness is only a matter of
formulation. The absence of shift invariance, in contrast, matters to the
placement of the coordinate sets.

\begin{lemma}[Pattern submultiplicativity and deletion]\label{lem:pattern-submult}
Let $\Omega\subset E^{\Nzero}$ be a pattern system and $q=\#E$.
For $m,n\in\N$ and finite $B\subset F\subset\Nzero$,
\begin{align}
 p_\Omega^*(m+n)&\le p_\Omega^*(m)p_\Omega^*(n),
 \label{eq:pattern-submult}\\
 \Patt_\Omega(F)&\le q^{\#(F\setminus B)}\Patt_\Omega(B).
 \label{eq:deletion} 
\end{align}
If $B\subset G$, then $\Patt_\Omega(B)\le\Patt_\Omega(G)$.
Consequently,
\begin{equation}\label{eq:fekete-pstar}
 \hpat(\Omega;\tau)
 =\inf_{n\ge1}\frac1{n\tau}\log p_\Omega^*(n),
 \qquad
 \log p_\Omega^*(n)\ge n\tau\hpat(\Omega;\tau).
\end{equation}
\end{lemma}

\begin{proof}
Split any $(m+n)$-element set into disjoint subsets of sizes $m$ and $n$.
Restriction sends a pattern on the union injectively to the pair of its two
restrictions, proving \eqref{eq:pattern-submult}.  Restriction from $F$ to
$B$ has fibers of cardinality at most $q^{\#(F\setminus B)}$, which proves
\eqref{eq:deletion}.  If $B\subset G$, every pattern on $B$ extends to a
pattern on $G$ using the same element of $\Omega$, giving monotonicity.
Fekete's lemma applied to $\log p_\Omega^*(n)$ proves
\eqref{eq:fekete-pstar}.
\end{proof}

For comparison, let $\Omega=\{0^\infty,10^\infty\}$. This set is a
closed one-sided subshift, but
$\Patt_\Omega(\{0\})=2$ and $\Patt_\Omega(\{m\})=1$ for $m\ge1$.
One cannot simply translate a maximizing set. Deleting a finite initial
interval has a different advantage: the number of lost coordinates is
bounded independently of the size of the maximizing set. Inequality
\eqref{eq:deletion} turns this bound into a controlled information loss.

\begin{theorem}[Simultaneous attainment]\label{thm:A}
Fix $\tau>0$.  For every countable family
$(\Omega_j)_{j\ge1}$ of finite-alphabet pattern systems, there is one
$A\in\mathscr S$ such that
\begin{equation}\label{eq:simultaneous-attainment}
 \hseq^A(\Omega_j;\tau)=\hpat(\Omega_j;\tau)
 \qquad(j\ge1).
\end{equation}
Thus the maximal-pattern entropy rate of each system is attained along the same
sequence.  No closedness, compactness, shift invariance, shift surjectivity, or
backward extension is required.
\end{theorem}

\begin{proof}
Let $E_j$ be the alphabet of $\Omega_j$, put $q_j=\#E_j$ and
$h_j^*=\hpat(\Omega_j;\tau)$.  For every $A$ and $n$,
\begin{align*}
 \Patt_{\Omega_j}(\{a_1,\ldots,a_n\})\le p_{\Omega_j}^*(n),
\end{align*}
so \eqref{eq:fekete-pstar} gives
$\hseq^A(\Omega_j;\tau)\le h_j^*$.

Choose an enumeration of stages
\begin{align*}
 (j_k)_{k\ge1}=(1;\ 1,2;\ 1,2,3;\ 1,2,3,4;\ldots),
\end{align*}
where semicolons only separate finite groups. Every positive integer
occurs infinitely often. Choose $\eta_k>0$ with $\eta_k\to0$.  We recursively
construct nonempty finite sets $B_k$, each lying strictly to the right of all
previous ones.  Put
\begin{align*}
 G_{k-1}=\bigcup_{i<k}B_i,
 \qquad r_{k-1}=\#G_{k-1},
 \qquad M_{k-1}=\max G_{k-1},
\end{align*}
with $r_0=0$ and $M_0=-1$.

If $h_{j_k}^*>0$, choose $n_k$ so large that
\begin{equation}\label{eq:explicit-losses}
 n_k>M_{k-1}+1,
 \qquad
 \frac{(M_{k-1}+1)\log q_{j_k}}{n_k\tau}<\eta_k,
 \qquad
 \frac{r_{k-1}}{n_k}<\eta_k.
\end{equation}
The quantities on the numerators are fixed before $n_k$ is chosen,
so these requirements can be satisfied even if the previously selected
coordinates are extremely far apart. Notice that $M_{k-1}+1$ bounds the
number of \emph{possible} old positions, whereas $r_{k-1}$ counts the
positions actually selected. The former controls deletion; the latter
controls dilution by historical observations. No uniform bound on $q_j$
or on the coordinate gaps is required. Choose an
$n_k$-element $F_k$ attaining $p_{\Omega_{j_k}}^*(n_k)$ and set
\begin{align*}
 B_k=F_k\cap\{M_{k-1}+1,M_{k-1}+2,\ldots\}.
\end{align*}
At most $M_{k-1}+1$ coordinates are deleted, and the first condition in
\eqref{eq:explicit-losses} makes $B_k$ nonempty.  If $h_{j_k}^*=0$, set
$B_k=\{M_{k-1}+1\}$.

Let $A$ be the increasing enumeration of $\bigcup_{k\ge1}B_k$.  Fix $j$
with $h_j^*>0$ and restrict to the infinitely many stages with $j_k=j$.
Write $b_k=\#B_k$ and $R_k=r_{k-1}+b_k=\#G_k$; then $G_k$ consists of the
first $R_k$ entries of $A$.  By deletion, \eqref{eq:fekete-pstar}, and
monotonicity,
\begin{align}
 \log\Patt_{\Omega_j}(G_k)
 &\ge \log\Patt_{\Omega_j}(B_k)\notag\\
 &\ge \log p_{\Omega_j}^*(n_k)-(M_{k-1}+1)\log q_j\notag\\
 &\ge n_k\tau h_j^*-(M_{k-1}+1)\log q_j.
 \label{eq:endpoint-numerator}
\end{align}
Also $R_k\le r_{k-1}+n_k$.  At all sufficiently late such stages the last
line of \eqref{eq:endpoint-numerator} is positive.  Dividing by $R_k\tau$
and using \eqref{eq:explicit-losses} yields the explicit endpoint estimate
\begin{equation}\label{eq:endpoint-bound}
 \frac1{R_k\tau}\log\Patt_{\Omega_j}(G_k)
 >\frac{h_j^*-\eta_k}{1+\eta_k}.
\end{equation}
For a fixed $j$, the indices $k$ with $j_k=j$ tend to infinity, hence
$\eta_k\to0$ along those stages. Since every $B_i$ is nonempty,
$R_k\ge k\to\infty$. The limsup along these
endpoints is at least $h_j^*$, proving the lower
bound.  If $h_j^*=0$, the universal upper bound and nonnegativity give
equality for every $A$.  This proves \eqref{eq:simultaneous-attainment} for
all $j$.
\end{proof}

\begin{remark}[What is attained]\label{rem:attainment}
The theorem asserts equality of limsup rates. It does not assert a common
subsequence of endpoints on which every system is simultaneously close
to its maximal rate. The good endpoints used for different $j$ can be
different. Nor does it assert convergence of the normalized counts.
For example, take $\tau=1$, let $B\subset\Nzero$ and its complement both be infinite,
and let $\Omega_B$ consist of all binary sequences supported on $B$.
For any schedule, the two normalized logarithmic counts for $\Omega_B$
and $\Omega_{\Nzero\setminus B}$ add to $\log2$ at every $n$, although
both maximal-pattern rates are $\log2$. They cannot both converge to
$\log2$. The two limsups, however, can both equal $\log2$.

Exact finite-level maximization is also a stronger and generally false
requirement. Put
\begin{align*}
 \Omega=\{0x:x\in\{0,1\}^{\Nzero}\}\cup\{10^\infty,20^\infty\}
 \subset\{0,1,2\}^{\Nzero}.
\end{align*}
The unique maximizing singleton is $\{0\}$, with three patterns.
For a three-element set containing zero there are six patterns, whereas
any three positive coordinates realize eight patterns. No increasing
sequence maximizes the counts at both levels $n=1$ and $n=3$.
\end{remark}

\begin{corollary}[Fixed-language attainment]\label{cor:fixed}
For every finite-alphabet pattern system $\Omega$ and every $\tau>0$,
\begin{equation}\label{eq:fixed-attainment}
 \hpat(\Omega;\tau)=\max_{A\in\mathscr S}\hseq^A(\Omega;\tau).
\end{equation}
The same identity holds for every finite Borel invariant partition $\C$ after
putting $\Omega=\Sigma_\C$.
\end{corollary}

\section{A language-quotient minimax theorem}\label{sec:minimax}

Fix a compact controlled-invariant set $Q$ and one period $\tau\in\N$.
For two invariant partitions $\C$ and $\D$ of this period, write
$\C\sim_{\lang}\D$ if there is a bijection
$\pi:I_\C\to I_\D$ such that, for every $N\in\N$,
\begin{equation}\label{eq:language-equivalence}
 \pi^N(W_N(\C;Q))=W_N(\D;Q),
\end{equation}
where $\pi^N$ acts coordinatewise.  This is \emph{language equivalence}.
Identity, inverse, and composition of the alphabet bijections show that
$\sim_{\lang}$ is an equivalence relation.  If $F\subset\Nzero$ is finite and
$m=\max F$, then \eqref{eq:language-equivalence} at length $m+1$, followed by
projection onto $F$, gives a bijection between the two $F$-pattern sets.
Thus language-equivalent partitions have identical finite-coordinate pattern
counts and identical sequence and maximal-pattern complexities.

For a nonempty class $\mathfrak C\subset\mathfrak C_{\Bor}(Q,\tau)$, set
\begin{align}
 H_{\rm pat}^*(Q;\mathfrak C)
 &=\inf_{\C\in\mathfrak C}\hpat(\C;Q),
 \label{eq:system-pattern}\\
 H_{\rm seq}^A(Q;\mathfrak C)
 &=\inf_{\C\in\mathfrak C}\hseq^A(\C;Q).
 \label{eq:system-sequence}
\end{align}
Every individual rate is nonnegative and finite, and because the class is nonempty both
infima are finite; no uniform bound on the numbers of atoms is needed.

The elementary minimax inequality gives only
\begin{align*}
 \sup_A\inf_{\C\in\mathfrak C}\hseq^A(\C;Q)
 \le\inf_{\C\in\mathfrak C}\sup_A\hseq^A(\C;Q)
 =H_{\rm pat}^*(Q;\mathfrak C).
\end{align*}
Separate maximizing schedules for separate partitions do not give the
opposite inequality. For instance, a two-by-two array with zero diagonal
and unit off-diagonal has maximum row minimum zero and minimum column
maximum one. The next theorem avoids this logical obstruction by
attaining every individual rate at one schedule.

\begin{theorem}[Minimax over countably many language types]\label{thm:B}
Let $\varnothing\ne\mathfrak C\subset\mathfrak C_{\Bor}(Q,\tau)$ and suppose
that the quotient $\mathfrak C/{\sim_{\lang}}$ is finite or countable.  There
is one $A\in\mathscr S$ such that
\begin{equation}\label{eq:pointwise-class-attainment}
 \hseq^A(\C;Q)=\hpat(\C;Q)
 \qquad(\C\in\mathfrak C).
\end{equation}
Consequently,
\begin{equation}\label{eq:minimax}
 \max_{B\in\mathscr S}\inf_{\C\in\mathfrak C}
 \hseq^B(\C;Q)
 =\inf_{\C\in\mathfrak C}\hpat(\C;Q).
\end{equation}
The outer maximum is attained; the partition infimum need not be attained.
\end{theorem}

\begin{proof}
Choose one representative $\C_j$ of each language-equivalence class.
Enumerate the representatives, repeating them periodically when the quotient
is finite, and apply Theorem~\ref{thm:A} to the resulting countable family
$\Sigma_{\C_j}$, using the common scale $\tau$.  The resulting $A$ gives equality for each
representative.  Language equivalence preserves all finite pattern counts, so
it gives \eqref{eq:pointwise-class-attainment} for every $\C\in\mathfrak C$.
Taking the infimum over $\C$ gives equality at this $A$.

For any $B\in\mathscr S$, Corollary~\ref{cor:fixed} gives
$\hseq^B(\C;Q)\le\hpat(\C;Q)$ for every $\C$.  Taking first the infimum over
$\C$ and then the supremum over $B$ proves the reverse minimax inequality and
hence \eqref{eq:minimax}.  If the pattern infimum is zero, the pointwise upper
bound and nonnegativity in fact give
$\inf_{\C\in\mathfrak C}\hseq^B(\C;Q)=0$ for every $B\in\mathscr S$, without
any countability assumption.
\end{proof}

The hypothesis concerns the symbolic behavior generated on $Q$, not the
names of controls used to present a strategy.  It identifies alphabet renumberings and strategies with the same
admissible language after such a renumbering. It does not merge distinct
atoms merely because their assigned controls agree.  It is therefore closer to an intrinsic
statement about the chosen partition class, although it remains relative to
that class and to the fixed period.

\begin{proposition}[A verifiable clopen sufficient condition]
\label{prop:digital}
Assume that $Q$ is compact metrizable and that the control alphabet $U$ is
countable.  Then
$\mathfrak C_{\clop}(Q,\tau)$ is countable.  Hence, if this class is nonempty,
Theorem~\ref{thm:B} applies to it.  In particular this holds when $U$ is a
finite digital alphabet.
\end{proposition}

\begin{proof}
A compact metrizable space has a countable base.  Every clopen subset is
compact and hence is a finite union of basic open sets contained in it.  Thus
$Q$ has only countably many clopen subsets and only countably many finite
clopen partitions.  Since $\tau$ is finite, $U^\tau$ is countable.  A finite
partition admits only countably many block assignments.  The invariant assignments form a
subcollection.
\end{proof}

Theorem~\ref{thm:B} fixes $\tau$.  It does not produce one common
physical-time schedule for partitions of different periods and proves no
minimax identity over $\mathfrak C_{\Bor}(Q)$.

\begin{example}[A purely combinatorial uncountable obstruction]
\label{ex:uncountable}
For every infinite $B\subset\Nzero$, let
\begin{align*}
 \Omega_B=\{x\in\{0,1\}^{\Nzero}:\supp(x)\subset B\}.
\end{align*}
At scale $\tau=1$,
\begin{align*}
 \Patt_{\Omega_B}(F)=2^{\#(F\cap B)},
 \qquad \hpat(\Omega_B;1)=\log2.
\end{align*}
For $A=(a_i)\in\mathscr S$,
\begin{align*}
 \hseq^A(\Omega_B;1)
 =\log2\cdot\limsup_{n\to\infty}
 \frac{\#(B\cap\{a_1,\ldots,a_n\})}{n}.
\end{align*}
If $\Nzero\setminus\{a_i:i\ge1\}$ is infinite, choose an infinite $B$ in that
complement.  Otherwise choose $B=\{a_{k^2}:k\ge1\}$.  In either case the
displayed upper density is zero.  Therefore
\begin{equation}\label{eq:abstract-obstruction}
 \sup_{A\in\mathscr S}
 \inf_{\substack{B\subset\Nzero\\B\text{ infinite}}}
 \hseq^A(\Omega_B;1)=0
 <\log2
 =\inf_{\substack{B\subset\Nzero\\B\text{ infinite}}}
 \hpat(\Omega_B;1).
\end{equation}
This proves only that minimax exchange is not a formal consequence of
pointwise attainment for arbitrary uncountable pattern families.  Every
invariant-partition name system satisfies
$\sigma(\Sigma_\C)\subset\Sigma_\C$ by suffix closure.  In contrast,
$\sigma(\Omega_B)\subset\Omega_B$ forces $k-1\in B$ whenever
$k\in B\setminus\{0\}$; an infinite such $B$ must equal $\Nzero$.  Hence the
proper sparse sets producing \eqref{eq:abstract-obstruction} cannot occur as
these autonomous control name systems.  The equation is not a deterministic
control counterexample and yields no failure statement for the unrestricted
Borel invariant-partition class.
\end{example}

\section{Realizing one-sided subshifts as safety systems}\label{sec:realization}

Let $E$ be a finite nonempty discrete alphabet and let
$\varnothing\ne Y\subset E^{\Nzero}$ be closed with $\sigma(Y)\subset Y$.
Write
\begin{align*}
 \Lcal_N(Y)=\{y_0\cdots y_{N-1}:y\in Y\},
 \qquad
 P_Y(F)=\#\{y|_F:y\in Y\}.
\end{align*}
The coordinate complexities are
\begin{align*}
 h_{\rm seq,coord}^A(Y)
 &=\limsup_{n\to\infty}\frac1n
   \log P_Y(\{a_1,\ldots,a_n\}),\\
 h_{\rm pat,coord}^*(Y)
 &=\lim_{n\to\infty}\frac1n
   \log\max_{\substack{F\subset\Nzero\\\#F=n}}P_Y(F).
\end{align*}

Adjoin an isolated cemetery point and give
$X=Y\sqcup\{\dagger\}$ the topological disjoint-union topology.  Put $U=E$
and define
\begin{equation}\label{eq:matching-control}
 f(y,u)=
 \begin{cases}
   \sigma y,&y_0=u,\\
   \dagger,&y_0\ne u,
 \end{cases}
 \qquad
 f(\dagger,u)=\dagger.
\end{equation}

\begin{theorem}[Symbolic-to-control realization]\label{thm:C}
For the matching-control system \eqref{eq:matching-control}:
\begin{enumerate}
\item $X$ is compact metrizable, $f:X\times U\to X$ is continuous, and $Y$
is a compact controlled-invariant set.
\item For every $N\in\N$,
\begin{equation}\label{eq:rinv-language}
 r_{\rm inv}(N,Y)=\#\Lcal_N(Y),
\end{equation}
and therefore
\begin{equation}\label{eq:hinv-htop}
 \hinv(Y)=\htop(Y,\sigma)
 =\lim_{N\to\infty}\frac1N\log\#\Lcal_N(Y).
\end{equation}
\item At period one identify $U^1$ with $U$.  If
$\C=(\{D_i:i\in I_\C\},1,\nu)\in\mathfrak C_{\Bor}(Y,1)$, set
$\lambda_\C(i)=\nu(D_i)\in E$.  Then every atom satisfies
\begin{equation}\label{eq:atom-cylinder}
 D_i\subset[\lambda_\C(i)]_Y
 :=\{y\in Y:y_0=\lambda_\C(i)\}.
\end{equation}
Moreover, the one-block label map
\begin{equation}\label{eq:label-factor}
 \Lambda_\C:\Sigma_\C\longrightarrow Y,
 \qquad (a_j)_{j\ge0}\longmapsto(\lambda_\C(a_j))_{j\ge0},
\end{equation}
is a continuous, onto, shift-commuting map.
\item For every $A\in\mathscr S$,
\begin{align}
 \inf_{\C\in\mathfrak C_{\Bor}(Y,1)}\hseq^A(\C;Y)
 &=\inf_{\C\in\mathfrak C_{\clop}(Y,1)}\hseq^A(\C;Y)
 =h_{\rm seq,coord}^A(Y),
 \label{eq:seq-realization}\\
 \inf_{\C\in\mathfrak C_{\Bor}(Y,1)}\hpat(\C;Y)
 &=\inf_{\C\in\mathfrak C_{\clop}(Y,1)}\hpat(\C;Y)
 =h_{\rm pat,coord}^*(Y),
 \label{eq:pat-realization}\\
 \inf_{\C\in\mathfrak C_{\Bor}(Y,1)}h(\C;Y)
 &=\inf_{\C\in\mathfrak C_{\clop}(Y,1)}h(\C;Y)
 =\htop(Y,\sigma).
 \label{eq:ordinary-realization}
\end{align}
Every Borel and clopen infimum in
\eqref{eq:seq-realization}--\eqref{eq:ordinary-realization} is attained by
the same first-coordinate cylinder partition.
\end{enumerate}
\end{theorem}

\begin{proof}
The space $X$ is the union of the two clopen compact sets $Y$ and
$\{\dagger\}$, hence is compact metrizable.  For fixed $u\in E$, the sets
$[u]_Y$, $Y\setminus[u]_Y$, and $\{\dagger\}$ are clopen; on them $f_u$ is,
respectively, the shift or a constant.  Thus $f_u$ is continuous.  Since $U$
is finite discrete, $f$ is jointly continuous.  Choosing $u=y_0$ sends $y$
to $\sigma y\in Y$, proving controlled invariance.

A direct induction gives
\begin{equation}\label{eq:safe-iff-match}
 \varphi(n,y,\omega)\in Y\ (0\le n\le N)
 \quad\Longleftrightarrow\quad
 \omega_j=y_j\ (0\le j<N).
\end{equation}
Indeed, after $j$ matches the state is $\sigma^j y$, whereas a mismatch enters
$\dagger$ permanently.  Hence a length-$N$ control word is safe from at
least one state of $Y$ exactly when it belongs to $\Lcal_N(Y)$, and then it
safely covers precisely its nonempty prefix cylinder.  Distinct language
words give disjoint cylinders. More explicitly:

\emph{Lower bound.} For each $w\in\Lcal_N(Y)$ choose $y^w\in Y$
with prefix $w$. If $S\subset U^N$ spans $Y$, some member of $S$ is safe
from $y^w$. By \eqref{eq:safe-iff-match} that member must equal $w$.
Thus $\Lcal_N(Y)\subset S$, and
$r_{\rm inv}(N,Y)\ge\#\Lcal_N(Y)$.

\emph{Upper bound.} Take $S=\Lcal_N(Y)$. For any $y\in Y$, its
prefix $y_0\cdots y_{N-1}$ belongs to $S$ and is safe through time $N$
by \eqref{eq:safe-iff-match}. Hence
$r_{\rm inv}(N,Y)\le\#\Lcal_N(Y)$.
This proves \eqref{eq:rinv-language}.  The language counts are
submultiplicative, so their normalized logarithms converge.

For completeness, this limit is the usual topological entropy even when
$\sigma|_Y$ is not onto.  The join of the first-coordinate cylinder cover
through times $0,\ldots,N-1$ has exactly $\#\Lcal_N(Y)$ nonempty atoms.
Conversely, given a finite open cover of compact $Y$, the partition into
cylinders of some fixed length $m$ refines that cover; its $N$-fold dynamical join is
then refined by the length-$(N+m-1)$ cylinders.  Taking exponential rates proves
\eqref{eq:hinv-htop}.

Let $\C=(\{D_i:i\in I_\C\},1,\nu)$ be a Borel invariant partition.  If
$D_i$ contained $y$ with $y_0\ne\lambda_\C(i)$, its assigned control would send $y$ to
$\dagger\notin Y$, contradicting safety.  This proves
\eqref{eq:atom-cylinder}.

Take $a=(a_j)_{j\ge0}\in\Sigma_\C$ and put $z_j=\lambda_\C(a_j)$.  For each
$N$, choose a witness $y^{(N)}$ for the first $N$ symbols of $a$.  Induction
using matching controls and \eqref{eq:atom-cylinder} gives
$y^{(N)}_j=\lambda_\C(a_j)=z_j$ for $0\le j<N$.  Hence
$y^{(N)}\to z$ in the product topology; since $Y$ is closed, $z\in Y$.  Thus
$\Lambda_\C$ is well defined, and it is continuous: the preimage of a cylinder in $E^{\Nzero}$ is a finite
union of cylinders in $I_\C^{\Nzero}$, restricted to $\Sigma_\C$.
It is shift commuting because
$(\Lambda_\C(\sigma a))_j=\lambda_\C(a_{j+1})
=(\sigma\Lambda_\C(a))_j$.  Given $y\in Y$, let $a_j$ be the unique index with
$\sigma^j y\in D_{a_j}$.  Its assigned label is $y_j$, and $y$ witnesses
every finite prefix of this atom itinerary.  Thus \eqref{eq:label-factor} is onto.  Here each prefix witness $y^{(N)}$ may be different. Its first $N$
\emph{control labels} nevertheless equal those of $z$, so convergence to
$z\in Y$ follows. This does not prove that $z$ follows the original atom
name $a$; its orbit may meet different atoms with the same labels.
Thus neither shift surjectivity nor a common atom-name witness is used.

For every finite $F$, the coordinate projection of $\Lambda_\C$ is onto, so
\begin{equation}\label{eq:partition-dominates-coordinate}
 \Patt_\C(F)\ge P_Y(F).
\end{equation}
It follows that every Borel, and hence every clopen, invariant partition has
sequence, maximal-pattern, and consecutive-name entropy at least the
corresponding coordinate quantity.

Let $E_Y=\{e\in E:[e]_Y\ne\varnothing\}$ and define
\begin{align*}
\C_0=(\{D_e=[e]_Y:e\in E_Y\},1,\nu_0),
 \qquad \nu_0(D_e)=e.
\end{align*}
This is a clopen invariant partition.  Under the symbol identification,
$W_N(\C_0;Y)=\Lcal_N(Y)$ and $\Sigma_{\C_0}=Y$; the latter follows from
closedness of $Y$.  Equality holds in \eqref{eq:partition-dominates-coordinate}
for every $F$.  Thus $\C_0$ attains all the lower bounds, proving
\eqref{eq:seq-realization}--\eqref{eq:ordinary-realization}.
\end{proof}

\section{A strict Thue--Morse separation}\label{sec:thue-morse}

Let $\vartheta(0)=01$, $\vartheta(1)=10$, and let
\begin{align*}
 \mathbf t=\lim_{r\to\infty}\vartheta^r(0)
 =0110100110010110\cdots.
\end{align*}
Equivalently, $t_n=s_2(n)\pmod2$, where $s_2(n)$ is the number of ones in
the binary expansion of $n$.  Let
\begin{equation}\label{eq:thue-morse-shift}
 Y=\overline{\{\sigma^m\mathbf t:m\ge0\}}
 \subset\{0,1\}^{\Nzero}.
\end{equation}
This is the one-sided Thue--Morse orbit closure.  Its full maximal pattern
complexity is classical \cite{KamaeZamboniWords2002,KamaeZamboniSystems2002}; the next
proof is included to make both rates used in the control corollary explicit.

\begin{proposition}[Exact symbolic rates]\label{prop:thue-morse}
For every $N\ge1$,
\begin{equation}\label{eq:TM-linear}
 \#\Lcal_N(Y)\le8N.
\end{equation}
Hence $\htop(Y)=0$.  On the single schedule
\begin{equation}\label{eq:TM-schedule}
 A_{\rm TM}=(a_i)_{i\ge1},
 \qquad a_i=4^{i-1}-1,
\end{equation}
one has, for every $k\ge1$,
\begin{equation}\label{eq:TM-full-pattern}
 P_Y(\{a_1,\ldots,a_k\})=2^k.
\end{equation}
The chosen coordinate set contains zero.  Since a cylinder is clopen, its
patterns in the orbit closure are exactly the patterns obtained along shifts
of $\mathbf t$; thus this calculation also matches the anchored word-pattern
convention of \cite{KamaeZamboniWords2002}.  Consequently
\begin{equation}\label{eq:TM-complexities}
 h_{\rm seq,coord}^{A_{\rm TM}}(Y)
 =h_{\rm pat,coord}^*(Y)=\log2.
\end{equation}
\end{proposition}

\begin{proof}
A finite word is in $\Lcal_N(Y)$ exactly when it occurs as a factor of
$\mathbf t$: one direction follows from the orbit, and the other from
convergence in a prefix cylinder.  Choose $\ell=\lceil\log_2N\rceil$.  The
identity $\mathbf t=\vartheta^\ell(\mathbf t)$ decomposes $\mathbf t$ into
aligned supertiles of length $2^\ell$.  Since $N\le2^\ell$, every length-$N$
factor is contained in $\vartheta^\ell(ab)$ for an adjacent two-letter factor
$ab$.  There are at most four choices for $ab$ and at most $2^\ell$ starting
offsets in the first supertile.  Therefore
$\#\Lcal_N(Y)\le4\,2^\ell\le8N$, proving zero topological entropy.

Fix $k\ge1$ and
$\varepsilon=(\varepsilon_0,\ldots,\varepsilon_{k-1})\in\{0,1\}^k$.
Put
\begin{align*}
 r=\#\{0\le i<k:\varepsilon_i=0\},
 \qquad
 q=\sum_{\substack{0\le i<k\\\varepsilon_i=0}}4^i
 +(r\bmod2)4^k.
\end{align*}
Only even binary positions occur in this sum.  The binary digit sum of $q$ is
$r+(r\bmod2)$, which is even, so $t_q=0$.  If $\varepsilon_i=1$, bit $2i$
of $q$ is zero, and adding $4^i=2^{2i}$ changes the digit-sum parity.  If
$\varepsilon_i=0$, bit $2i$ is one and bit $2i+1$ is zero; adding $2^{2i}$
replaces these two bits by zero and one without a further carry, so the parity
is unchanged.  In both cases
\begin{equation}\label{eq:TM-digit-identity}
 t_{q+4^i}=\varepsilon_i\qquad(0\le i<k).
\end{equation}
The extra bit at position $2k$ only corrects the initial parity of $q$.
For $i<k$, a carry stops at the empty odd position $2i+1$, which is
strictly below $2k$. The parity correction therefore interferes with
none of these tests. Each addition $q+4^i$ is performed separately;
we are not adding all the $4^i$'s successively to one integer.

Set $n=q+1$. Since $a_{i+1}=4^i-1$,
\begin{align*}
 (\sigma^n\mathbf t)_{a_{i+1}}
 =t_{q+4^i}=\varepsilon_i.
\end{align*}
Thus every binary $k$-pattern occurs on the first $k$ members of
$A_{\rm TM}$, proving \eqref{eq:TM-full-pattern}.  The binary alphabet gives
the reverse upper bound $p_Y^*(k)\le2^k$, and
\eqref{eq:TM-complexities} follows.
\end{proof}

\begin{corollary}[Zero invariance entropy, full binary sparse-time complexity]
\label{cor:separation}
Apply the matching-control construction \eqref{eq:matching-control} to the
Thue--Morse subshift \eqref{eq:thue-morse-shift}.  Then
\begin{equation}\label{eq:strict-control-separation}
 \hinv(Y)=0
 <\log2
 =\inf_{\C\in\mathfrak C_{\Bor}(Y,1)}\hpat(\C;Y)
 =\inf_{\C\in\mathfrak C_{\clop}(Y,1)}\hpat(\C;Y).
\end{equation}
Moreover, the explicit schedule $A_{\rm TM}$ satisfies
\begin{align}
 \max_{A\in\mathscr S}\inf_{\C\in\mathfrak C_{\Bor}(Y,1)}
 \hseq^A(\C;Y)
 &=\inf_{\C\in\mathfrak C_{\Bor}(Y,1)}\hseq^{A_{\rm TM}}(\C;Y)
 =\log2,
 \label{eq:TM-Borel-sequence}\\
 \max_{A\in\mathscr S}\inf_{\C\in\mathfrak C_{\clop}(Y,1)}
 \hseq^A(\C;Y)
 &=\inf_{\C\in\mathfrak C_{\clop}(Y,1)}\hseq^{A_{\rm TM}}(\C;Y)
 =\log2.
 \label{eq:TM-clopen-sequence}
\end{align}
\end{corollary}

\begin{proof}
Combine Theorem~\ref{thm:C} with Proposition~\ref{prop:thue-morse}.  For any
schedule the binary coordinate complexity is at most $\log2$, while
$A_{\rm TM}$ attains $\log2$.  Equations
\eqref{eq:seq-realization}--\eqref{eq:pat-realization} give the Borel and
clopen infima.
\end{proof}

Corollary~\ref{cor:separation} is a deterministic control-system statement,
not merely a symbolic analogy.  It also shows that the new sparse-time
invariant-language infimum cannot in general be identified with Kawan's
invariance entropy.

\section{Comparison, all periods, and limitations}\label{sec:comparison}

\subsection{Comparison with Kawan's invariance entropy}
For every invariant partition, consecutive coordinates are one permissible
choice in the maximal-pattern count. Therefore
\begin{equation}\label{eq:ordinary-pattern-one-sided}
 h(\C;Q)\le\hpat(\C;Q)\le\frac{\log\#I_\C}{\tau}.
\end{equation}
Kawan's Theorem~2.3, with natural logarithms, gives
\begin{equation}\label{eq:kawan-characterization}
 \hinv(Q)=\inf_{\C\in\mathfrak C_{\Bor}(Q)}h(\C;Q).
\end{equation}
Its hypotheses are satisfied by the full-control discrete-time systems
used here: $X$ is metrizable, each controlled map is continuous, and
$Q$ is nonempty, compact and controlled invariant.
The period in this formula is allowed to vary.

For clarity, the relevant argument can be given directly. Spanning
families concatenate, so
$r_{\rm inv}(m+n,Q)\le r_{\rm inv}(m,Q)r_{\rm inv}(n,Q)$ when the
right-hand side is finite. Moreover $r_{\rm inv}(n,Q)$ is nondecreasing.
If any finite-horizon spanning number is finite, then $r_{\rm inv}(1,Q)$
is finite, all of them are finite by concatenation, and Fekete's lemma
gives a limit in \eqref{eq:hinv-definition}. If none is finite, that
limit equals $+\infty$.

Let $S_N=\{\omega^1,\ldots,\omega^m\}$ be a minimal finite
$N$-step spanning family. Its safe sets are
\begin{align*}
 K_i=Q\cap\bigcap_{s=1}^{N}
       \{x\in X:\varphi(s,x,\omega^i)\in Q\}.
\end{align*}
They are closed in $Q$, since finite control compositions are continuous
and $Q$ is closed. They cover $Q$. Set
$D_1=K_1$ and $D_i=K_i\setminus\bigcup_{j<i}K_j$, discard empty sets,
and assign $\omega^i$ to $D_i$. This gives a Borel invariant partition
$\C_N$ of period $N$ with at most $m$ atoms. This is the measurable
construction underlying Kawan's Theorem~2.3; Lemma~2.3 alone only
provides disjointization. The safe sets are already closed before
that lemma is applied.

For the other inequality in \eqref{eq:kawan-characterization}, all
admissible $k$-block words of any $\C$ form a spanning family after their
assigned controls are concatenated. Indeed, from each initial state the
feedback rule selects one such admissible word. Consequently
$r_{\rm inv}(k\tau,Q)\le\#W_k(\C;Q)$, and taking limits gives
$\hinv(Q)\le h(\C;Q)$. The reverse inequality follows from
$h(\C_N;Q)\le N^{-1}\log r_{\rm inv}(N,Q)$ and $N\to\infty$.
If no finite spanning family exists, no finite invariant partition exists,
and both sides are $+\infty$.

Thus, for every nonempty fixed-period subclass
$\mathfrak C\subset\mathfrak C_{\Bor}(Q,\tau)$, the general comparison is
\begin{equation}\label{eq:three-quantities}
 \hinv(Q)
 \le\inf_{\C\in\mathfrak C}h(\C;Q)
 \le\inf_{\C\in\mathfrak C}\hpat(\C;Q).
\end{equation}
The minimax theorem concerns the last quantity. The Thue--Morse
realization proves that it can strictly exceed $\hinv(Q)$.

\subsection{Why the unrestricted all-period infimum collapses}
There is a natural candidate for an all-period Borel quantity,
\begin{align*}
 H_{\rm pat,all}^{*,\Bor}(Q)
 :=\inf_{\tau\in\N}\inf_{\C\in\mathfrak C_{\Bor}(Q,\tau)}
       \hpat(\C;Q).
\end{align*}
The following proposition records the consequence of Kawan's spanning
construction for this all-period quantity.

\begin{proposition}[All-period collapse]\label{prop:all-period-collapse}
In the full-control discrete-time setting of Section~\ref{sec:control}, every
nonempty compact controlled-invariant set $Q$ satisfies
\begin{equation}\label{eq:all-period-collapse}
 H_{\rm pat,all}^{*,\Bor}(Q)=\hinv(Q).
\end{equation}
\end{proposition}

\begin{proof}
The lower bound follows from
\eqref{eq:three-quantities}. For the upper bound use the same
$\C_N$ just constructed. Since its alphabet has at most
$r_{\rm inv}(N,Q)$ symbols, every pattern on $k$ coordinates has at most
$r_{\rm inv}(N,Q)^k$ possibilities. Hence
\begin{align*}
 H_{\rm pat,all}^{*,\Bor}(Q)
 \le\hpat(\C_N;Q)
 \le\frac{\log r_{\rm inv}(N,Q)}{N}.
\end{align*}
Let $N\to\infty$. The infinite-entropy case was already covered by
nonexistence of finite invariant partitions. This proves
\eqref{eq:all-period-collapse}.
\end{proof}

For the matching-control Thue--Morse example, the collapse is particularly
visible. At period $N$, the length-$N$ cylinder partition with its matching
control blocks is clopen, has at most $8N$ atoms and has pattern rate at
most $N^{-1}\log(8N)$. Thus even the all-period clopen infimum is zero
in this example, while the period-one Borel infimum is $\log2$.
Changing the sampling period changes which information is packed into
one observed symbol and the cost assigned to that symbol.

\subsection{Countability, trajectories, and physical time}
The countability hypothesis in Theorem~\ref{thm:B} is used exactly when
one representative language from each class is listed and assigned
infinitely many construction stages. Compact metrizability gives
countably many clopen sets; it does not give countably many Borel sets
or Borel name languages. No approximation theorem allowing that
substitution is proved here. Example~\ref{ex:uncountable} rules out a
formal uncountable extension for arbitrary pattern sets, but its
non-shift-invariant examples do not settle the unrestricted autonomous
control problem. Theorem~\ref{thm:C} bypasses countability in its specific
realization by exhibiting a least language under the label projections.
Whether a useful general Borel minimax theorem can be established beyond
countably many language types remains outside the present results.

The patterns count admissible language projections. Every finite pattern
has a finite trajectory witness which stays safe between observations.
An infinite point in the completed Borel name system may have only
separate finite-prefix witnesses, as Remark~\ref{rem:witness} shows.
This distinction cannot be removed by compactness alone when the atoms
are not closed. The factor in Theorem~\ref{thm:C} is a factor of compact
\emph{name systems}; it is not a continuous factor defined by an arbitrary
Borel feedback map on the original state space.

For one fixed $\tau$, $A$ denotes block indices and $\tau A$ denotes
physical observation times. The normalization $n\tau$ charges one block
length per observation and disregards the lengths of the gaps. A common
block-index sequence for different periods would represent different
physical schedules. Theorem~\ref{thm:A} could be applied to countably many
languages with different positive normalizing constants by rescaling
each equality, but that would not solve the distinct physical-time
minimax problem. Formula~\eqref{eq:all-period-collapse} is an infimum
identity and makes no common-schedule assertion.

Finally, if one safe control block works for every point of $Q$, the
one-atom invariant partition makes all three fixed-period infima zero
for any class containing it. On a connected $Q$, every finite clopen
partition has one atom. The full clopen class is therefore either empty
or has zero pattern infimum, while Borel strategies can remain
nontrivial. These observations explain both the usefulness and the
limits of the clopen sufficient condition. The nontrivial separation
studied here records the richness of safe symbolic instructions at a
specified sampling period, despite a zero asymptotic count per elapsed
time.

\begin{thebibliography}{99}

\bibitem{Goodman1974}
T.~N.~T. Goodman,
Topological sequence entropy,
\emph{Proc. London Math. Soc.} (3) \textbf{29} (1974), 331--350,
\doi{10.1112/plms/s3-29.2.331}.

\bibitem{HuangMaassYe2004}
W.~Huang, A.~Maass and X.~Ye,
Sequence entropy pairs and complexity pairs for a measure,
\emph{Ann. Inst. Fourier (Grenoble)} \textbf{54} (2004), 1005--1028,
\doi{10.5802/aif.2041}.

\bibitem{HuangYe2009}
W.~Huang and X.~Ye,
Combinatorial lemmas and applications to dynamics,
\emph{Adv. Math.} \textbf{220} (2009), 1689--1716,
\doi{10.1016/j.aim.2008.11.009}.

\bibitem{Hulse1982}
P.~Hulse,
Sequence entropy and subsequence generators,
\emph{J. London Math. Soc.} (2) \textbf{26} (1982), 441--450,
\doi{10.1112/jlms/s2-26.3.441}.

\bibitem{KamaeUniform2009}
T.~Kamae,
Uniform sets and complexity,
\emph{Discrete Math.} \textbf{309} (2009), 3738--3747,
\doi{10.1016/j.disc.2008.10.001}.

\bibitem{KamaeZamboniWords2002}
T.~Kamae and L.~Zamboni,
Sequence entropy and the maximal pattern complexity of infinite words,
\emph{Ergodic Theory Dynam. Systems} \textbf{22} (2002), 1191--1199,
\doi{10.1017/S014338570200055X}.

\bibitem{KamaeZamboniSystems2002}
T.~Kamae and L.~Zamboni,
Maximal pattern complexity for discrete systems,
\emph{Ergodic Theory Dynam. Systems} \textbf{22} (2002), 1201--1214,
\doi{10.1017/S0143385702000585}.

\bibitem{Kawan2013}
C.~Kawan,
\emph{Invariance Entropy for Deterministic Control Systems: An Introduction},
Lecture Notes in Mathematics, vol.~2089, Springer, Cham, 2013,
\doi{10.1007/978-3-319-01288-9}.

\bibitem{TomarKawanZamani2022}
M.~S. Tomar, C.~Kawan and M.~Zamani,
Numerical over-approximation of invariance entropy via finite abstractions,
\emph{Systems Control Lett.} \textbf{170} (2022), article~105395,
\doi{10.1016/j.sysconle.2022.105395}.

\bibitem{Yao2026}
S.~Yao,
Invariance entropy in the dust,
arXiv:2607.02279v1 (2026),
\doi{10.48550/arXiv.2607.02279}.

\end{thebibliography}

\end{document}
